
\chapter{Операции и соединения}
\label{chap:connections}

В предыдущей главе мы создали лампу, которая
самостоятельно обрабатывала команды. Теперь научимся
связывать несколько автоматов в одну систему.

Для этого рассмотрим пример Array Sum --- пошагового 
вычисления суммы элементов массива.

Его логика разделена между двумя автоматами:

\begin{itemize}
  \item \code{SumController} выполняет вычисления,
    обрабатывая по одному элементу массива за шаг.
  \item \code{SumView} принимает действия пользователя
    и отвечает за отображение результатов.
\end{itemize}

Приложение поддерживает два режима: пошаговый, когда
пользователь управляет вычислением, и автоматический,
когда Controller самостоятельно выполняет все шаги.

Для обмена командами автоматы объявляют операции
и соединения между ними.

\section{Виды операций}

В CIAO операции делятся на предоставляемые
(\emph{provided}) и требуемые (\emph{required}).

Предоставляемая операция определяет, какие команды
или запросы автомат может принимать.
Требуемая --- какие операции других автоматов
ему необходимы.

Всего предусмотрено четыре вида операций:

\begin{table}[ht]
  \centering
  \caption{Виды операций CIAO}
  \begin{tabularx}{\textwidth}{@{}p{0.38\textwidth}X@{}}
    \toprule
    Декоратор & Назначение \\
    \midrule
    \code{@provided_command} &
      Принимает команду. \\
    \code{@required_command} &
      Отправляет команду другому автомату. \\
    \code{@provided_query} &
      Предоставляет ответ на запрос. \\
    \code{@required_query} &
      Запрашивает значение у другого автомата. \\
    \bottomrule
  \end{tabularx}
\end{table}

Команды передаются через очередь сообщений и
обрабатываются последовательно. Запросы выполняются
синхронно и сразу возвращают результат.
Их использование рассмотрим в
\cref{chap:guards-queries}.

\section{Соединения автоматов}

Для взаимодействия автоматов используется
\code{Connection}. Каждое соединение связывает
требуемую операцию одного объекта с предоставляемой
операцией другого.

Рассмотрим сборку программы Array Sum:

\lstinputlisting[
  caption={Сборка программы Array Sum},
  label={lst:array-sum-program}
]{code/array_sum_program.py}

Каждый конец соединения задаётся парой:
автоматный объект и операция его класса.

Например, соединение требуемой операции
\code{SumView.start} с предоставляемой
\code{SumController.start} позволяет View отправлять
команду запуска Controller.

\noindent\begin{minipage}{\textwidth}
  \centering
  \begin{tikzpicture}[
  automaton/.style={
    draw=ciaoBlue,
    fill=ciaoPaleBlue!35,
    very thick,
    rounded corners=2mm,
    minimum width=42mm,
    minimum height=23mm,
    align=center,
    font=\sffamily
  },
  connector/.style={draw=ciaoNavy,thick},
  operation/.style={font=\ttfamily\small,align=center}
]
  \node[automaton] (view) {
    {\scriptsize\color{ciaoMuted}\textless\textless ViewAutomaton\textgreater\textgreater}\\[-1mm]
    {\bfseries SumView}
  };
  \node[automaton,right=68mm of view] (controller) {
    {\scriptsize\color{ciaoMuted}\textless\textless ControllerAutomaton\textgreater\textgreater}\\[-1mm]
    {\bfseries SumController}
  };

  % Обе половины Connection завершаются в общей точке, как при экспорте
  % в draw.io: socket required-интерфейса принимает ball provided-интерфейса.
  \coordinate (junction) at ($(view.east)!0.5!(controller.west)$);

  % Required-половина от SumView до общего стыка.
  \draw[connector] (view.east) -- ($(junction)+(-3mm,0)$);
  \draw[connector]
    ($(junction)+(0,3mm)$)
    arc[start angle=90,end angle=270,radius=3mm];

  % Provided-половина от SumController; ball касается раскрытия socket.
  \coordinate (ball-center) at ($(junction)+(2.4mm,0)$);
  \draw[connector]
    (controller.west) -- ($(ball-center)+(2.4mm,0)$);
  \draw[connector,fill=white] (ball-center) circle[radius=2.4mm];

  % Экспорт draw.io подписывает всю assembly-связь один раз у её середины.
  \node[operation,above=6mm of junction] {start};
\end{tikzpicture}

  \captionof{figure}{Соединение required-команды View с provided-командой Controller}
  \label{fig:sum-start-connection}

  \vspace{0.4em}  
\end{minipage}

Первым аргументом \code{Connection} указывается
требуемая операция, вторым --- предоставляемая.

На этапе компиляции модели CIAO проверяются
направления операций, совместимость их сигнатур
и корректность соединений. Эта проверка выполняется
перед запуском Runtime.

Все объекты и соединения объединяются в \code{Program}.
Сборку программы удобно выносить в отдельную функцию
\code{build_program()}, отделяя создание системы
от её исполнения.

\section{Отправка команд}

Соединение определяет маршрут команды, но не
инициирует её отправку.

Для отправки команды автомат объявляет требуемую
операцию с помощью \code{@required_command}.
Её вызов выполняется внутри действия, отмеченного
декоратором \code{@action(emit=...)}.

Например:

\begin{lstlisting}[language=Python,numbers=none]
@action(emit=draw)
def publish(
    self, ctx: ControllerContext, emit: Emit
) -> None:
    emit(total=ctx.total)
\end{lstlisting}

Аргумент \code{emit=draw} объявляет требуемую
команду, которую может отправить действие.

При таком объявлении метод обязательно должен
принимать параметр типа \code{Emit}.
Это проверяется при компиляции модели CIAO.

Во время исполнения Runtime передаёт в этот
параметр объект для отправки команды.
Вызов \code{emit(...)} передаёт её аргументы,
а получатель определяется по \code{Connection}.

Если действие может отправлять несколько разных
команд, используется \code{@action(emits=(...))}
и параметр типа \code{Emits}.

Например, объявление
\code{@action(emits=(draw, continue_run))}
разрешает отправлять обе команды. Конкретная
операция выбирается при вызове:

\begin{lstlisting}[language=Python,numbers=none]
emits[self.draw](total=ctx.total)
emits[self.continue_run]()
\end{lstlisting}

Для \code{Emits} также обязательна соответствующая
аннотация параметра, проверяемая компилятором CIAO.

\begin{importantbox}
Рекомендуется использовать действия с \code{emit}
только для отправки команд, не изменяя в них контекст.

Если переход должен изменить контекст и отправить
команду, объявите отдельные Action и укажите их
в \code{actions} в нужном порядке.

На диаграмме CIAO вместо имени такого действия
отображается имя отправляемой required-операции.
\end{importantbox}

\begin{notebox}
Порядок выполнения Action определяется их расположением
в \code{actions}. Однако последовательность вызовов
\code{emit} внутри одного Action не выражена
в декларативной модели.

Аргумент \code{emits=(...)} перечисляет допустимые
команды, но не определяет порядок их отправки.
Он зависит от реализации Python-метода.
\end{notebox}

Отправленные команды помещаются в очередь после
успешного завершения перехода и обрабатываются
Runtime отдельными шагами.

В Array Sum этот механизм используется и для
автоматического вычисления: Controller соединяет
собственную требуемую команду \code{continue_run}
с предоставляемой \code{tick}.

При получении \code{tick} Controller обрабатывает один
элемент массива. В автоматическом режиме он также
отправляет \code{continue_run}, планируя обработку
следующего элемента.

Так вычисление продолжается до конца массива без
рекурсивных вызовов и участия пользовательского интерфейса.

\begin{notebox}
Не путайте команду \code{tick} с методом
\code{Runtime.step()}.

\code{tick} --- обычная команда автомата, название
и поведение которой определяет разработчик.
\code{Runtime.step()} --- метод исполнителя,
обрабатывающий одно сообщение из очереди.

Поэтому один вызов \code{Runtime.step()} не обязательно
соответствует одному шагу вычисления: в очереди могут
находиться и команды других автоматов.
\end{notebox}

\section{Внешние операции}

Не все операции предназначены для взаимодействия
между автоматами. Некоторые связывают программу
с внешним кодом, например с графическим интерфейсом.

Для этого используется параметр \code{external=True}:

\begin{itemize}
  \item У \code{@provided_command} он разрешает
    отправлять команду извне программы, например
    через \code{Runtime.send}.
  \item У \code{@required_command} он позволяет
    подключить внешний обработчик, например
    метод окна приложения.
\end{itemize}

Внутренние соединения между автоматами по-прежнему
задаются через \code{Connection}.

Таким образом, графический интерфейс может передавать
команды автоматам и получать от них результаты
без прямого доступа к их контекстам.

Работу с пользовательским интерфейсом рассмотрим
в \cref{chap:user-interface}.
